% !TeX root = ../../Book4.tex
% 纲要标题：### B.3 高阶代数 K 理论
% 纲要位置：工作记录/计划纲要.md，第 3882 行。
% 纲要细目（写作范围）：
% #### B.3.1 稳定一般线性群与一阶 K 群
% #### B.3.2 Q 构造与高阶 K 群
% #### B.3.3 局部化与高阶信息

\section{高阶代数 K 理论}
\label{b4:sec:B03}

\(K_0\) 将对象按可加关系组合，却没有直接保留对象的自同构。矩阵可逆变换提供了最初的高阶信息。再向上，关系之间的相容性需要用空间及其同伦群组织；这一节先把 \(K_1\) 算清，再说明一般 \(K_n\) 所依赖的准确构造。

\subsection{稳定一般线性群与一阶 K 群}
\label{b4:subsec:B03:01}

\begin{definition}\label{b4:def:B-k1}
设 \(R\) 为含幺环。通过 \(A\mapsto\operatorname{diag}(A,1)\) 把 \(\operatorname{GL}_n(R)\) 包含进 \(\operatorname{GL}_{n+1}(R)\)，所得并记为 \(\operatorname{GL}(R)\)，称为稳定一般线性群。其交换化
\[
K_1(R)\coloneqq
\operatorname{GL}(R)/\bigl[\operatorname{GL}(R),\operatorname{GL}(R)\bigr]
\]
称为 \(R\) 的\term[yijiedaishuKqun]{一阶代数 \(K\) 群}[algebraic \(K_1\)-group]。
\end{definition}
\symidx{\(\operatorname{GL}(R)\)：稳定一般线性群}
\symidx{\(K_1(R)\)：环的一阶代数 K 群}

交换子约定为 \([g,h]=ghg^{-1}h^{-1}\)。稳定化允许在矩阵右下角添加恒等块，而且确实可能消去固定大小交换化中的非零类。例如 \(\operatorname{GL}_2(\mathbb F_2)\) 的交换化为 \(C_2\)，而 \(K_1(\mathbb F_2)=\{1\}\)；其中二阶初等矩阵的非平凡类，在添加第三个指标后成为交换子的类。具体计算见\cref{b4:prop:B-stabilization-example}。定义可对照 Weibel\cite[Chapter III, Definition 1.1]{Weibel2013KBook}。

\begin{lemma}[稳定初等矩阵与交换子]\label{b4:lem:B-whitehead}
设 \(R\) 为含幺环，\(E(R)\) 为全部初等矩阵
\(e_{ij}(a)=I+aE_{ij}\)、\(i\ne j\)、\(a\in R\)
在 \(\operatorname{GL}(R)\) 中生成的子群。则
\[
E(R)=\bigl[\operatorname{GL}(R),\operatorname{GL}(R)\bigr].
\]
\end{lemma}
\begin{proof}
稳定化后总可选第三个指标 \(k\)。矩阵单位乘法给出
\(\bigl[e_{ik}(a),e_{kj}(1)\bigr]=e_{ij}(a)\)，故每个初等矩阵属于交换子群。

反向包含的目标是把任意交换子稳定后写成初等矩阵的乘积。为此，先把 \(g\) 与 \(g^{-1}\) 放在两个对角块上，证明 \(D(g)=\operatorname{diag}(g,g^{-1})\) 属于 \(E(R)\)；再用这些矩阵拼出交换子。要得到这样一对对角块，可以将带系数的交换矩阵与普通交换矩阵相乘，因此先对任意 \(g\in\operatorname{GL}_n(R)\) 写出
\[
w(g)=
\begin{pmatrix}I&g\\0&I\end{pmatrix}
\begin{pmatrix}I&0\\-g^{-1}&I\end{pmatrix}
\begin{pmatrix}I&g\\0&I\end{pmatrix}
=
\begin{pmatrix}0&g\\-g^{-1}&0\end{pmatrix}.
\]
两个形如上三角或下三角的块幺矩阵都是通常初等矩阵的乘积：在一个非对角块中的各矩阵单位两两乘积为零，逐个加入其系数即可。因此
\(D(g)=w(g)w(-I)=\operatorname{diag}(g,g^{-1})\)
属于 \(E(R)\)。对同样大小的 \(g,h\)，有
\[
D(g)D(h)D\bigl((hg)^{-1}\bigr)
=\operatorname{diag}\bigl([g,h],I\bigr).
\]
右下块为 \(g^{-1}h^{-1}hg=I\)，左上块为 \(ghg^{-1}h^{-1}\)。故每个交换子稳定后属于 \(E(R)\)，得到反向包含。
\end{proof}

这个证明给出 \(K_1(R)=\operatorname{GL}(R)/E(R)\)，并同时证明 \(E(R)\) 正规。无需事先断言每个固定大小的 \(E_n(R)\) 都在 \(\operatorname{GL}_n(R)\) 中正规；后一个断言对一般环并不成立。

\begin{theorem}\label{b4:thm:B-k1-field-integers}
对任意域 \(k\)，行列式诱导同构 \(K_1(k)\simeq k^\times\)。对整数环，有 \(K_1(\mathbb Z)\simeq\{1,-1\}\)。
\end{theorem}
\begin{proof}
行列式在稳定化下不变，且初等矩阵行列式为 \(1\)，故给出满同态 \(K_1(R)\rightarrow R^\times\)，其中 \(R=k\) 或 \(\mathbb Z\)。只需证明行列式为 \(1\) 的矩阵属于 \(E(R)\)。

域上可用行加法和带符号的行交换进行消元。带符号的交换矩阵
\(\left(\begin{smallmatrix}0&1\\-1&0\end{smallmatrix}\right)=w(1)\)
已由前一引理写成初等乘积。消元先得到非零对角元的上三角矩阵，再用行加法清除非对角项，得到行列式为 \(1\) 的对角矩阵。每个块 \(\operatorname{diag}(u,u^{-1})=D(u)\) 为初等乘积，因此任意乘积为 \(1\) 的对角元也可相继消去。

整数情形中，行列式为 \(1\) 使第一列元素的最大公因数为 \(1\)。两行间的带余除法与带符号交换，就是 Euclid 算法\footnote{欧几里得（\OriginalName{Εὐκλείδης}，约前325-约前265），古希腊数学家，著有《几何原本》，系统整理了几何学与数论。}，可把这一列化为 \((1,0,\ldots,0)^t\)；若先得到 \(-1\)，用 \(D(-1)\) 同时改变两行符号。然后用列加法清掉第一行剩余元素。余下方块仍为行列式 \(1\) 的整数矩阵，按大小归纳得到恒等矩阵。列加法也是右乘初等矩阵，故原矩阵仍在 \(E(\mathbb Z)\) 中。大小为 \(1\) 的情形直接成立；需要两行时可先稳定化一次。两种情形的行列式核均为初等子群，结合\cref{b4:lem:B-whitehead}完成证明。
\end{proof}

域上的 \(K_1(k)=k^\times\) 保留了乘法群的信息，而 \(K_0(k)=\mathbb Z\) 只反映维数。一般交换环还可能有非零的行列式核
\(\operatorname{SK}_1(R)=\ker\bigl(K_1(R)\rightarrow R^\times\bigr)\)，所以不能把域上的公式无条件推广。非零核的例子见 Weibel\cite[Chapter III, Example 1.5.4]{Weibel2013KBook}。

\subsection{Q 构造与高阶 K 群}
\label{b4:subsec:B03:02}

设 \(\mathcal E\) 为小正合范畴。Quillen 的 \(Q\) 构造保留对象，却以子商数据作为态射：从 \(A\) 到 \(B\) 的一张态射，指定 \(B\) 的一个可容许子对象，以及该子对象的一个可容许商 \(A\)。这不是原范畴中的一张 \(A\to B\) 映射。
\begin{definition}\label{b4:def:B-q-construction}
设 \(\mathcal E\) 为小正合范畴。构造范畴 \(Q\mathcal E\)，其对象与 \(\mathcal E\) 相同，从 \(A\) 到 \(B\) 的态射为图式
\[
A\xleftarrow{p}M\xrightarrow{i}B
\]
的同构类，其中 \(p\) 为可容许满态射，\(i\) 为可容许单态射。中间对象的同构必须使两端映射均相容。与 \(B\xleftarrow{q}N\xrightarrow{j}C\) 的复合由拉回 \(M\times_BN\) 给出：
\[
A\longleftarrow M\times_BN\longrightarrow C.
\]
恒等态射由两条恒等映射 \(A\leftarrow A\to A\) 表示。这个构造称为 Quillen 的 \(Q\) 构造。
\end{definition}
\symidx{\(Q\mathcal E\)：正合范畴的 Quillen \(Q\) 构造}

在模或向量空间中，用 \(i\) 把 \(M\) 识别为 \(B\) 的子对象，则复合中的拉回就是 \(q^{-1}(M)\subseteq N\)。我们先在 \(C\) 的子对象 \(N\) 中找出映入 \(M\) 的部分，再经 \(q\) 和 \(p\) 将它映到 \(A\)。因此新图式的左端仍是一个商，右端仍是一个子对象。特别地，从 \(0\) 到向量空间 \(B\) 的 \(Q\)-态射恰好对应 \(B\) 的子空间；当 \(B=\mathbb F_q^2\) 时共有 \(q+3\) 张，而非唯一一张。这个对应及计数见\cref{b4:prop:B-q-zero-subspaces}。

\ProseParagraphBreak
正合公理保证这个拉回存在，投影到 \(M\) 为可容许满态射。另取 \(i:M\rightarrow B\) 的可容许余核 \(B\rightarrow D\)，则复合 \(N\xrightarrow{q}B\rightarrow D\) 仍为可容许满态射。其核正是 \(M\times_BN\rightarrow N\)：一个到 \(N\) 的态射在复合到 \(D\) 后为零，当且仅当其到 \(B\) 的像经由 \(M\)，再由拉回泛性质得到唯一分解。因此该投影为可容许单态射，两端复合仍属所定类型。重复拉回的泛性质给出典范同构，通过中间对象同构类取商后，复合严格结合。恒等态射由 \(A\leftarrow A\rightarrow A\) 两个恒等映射表示。

\ProseParagraphBreak
范畴的神经 \(N\mathcal C\) 把可复合的 \(n\) 个态射作为 \(n\)-单形；删去端点或复合相邻态射给出面映射，插入恒等给出退化映射。其几何实现记为 \(B\mathcal C\)。这个空间同时反映对象、态射及复合的相容关系。本节假定读者熟悉同伦群等代数拓扑基础，不在附录中重述这些理论。
\symidx{\(N\mathcal C,\ B\mathcal C\)：范畴的神经与分类空间}

\ProseParagraphBreak
在 \(Q\mathcal E\) 中，每个对象 \(A\) 都有从 \(0\) 来的态射，所以 \(BQ\mathcal E\) 连通；\(\pi_0\) 无法区分对象的可加类。不过，从 \(0\) 到 \(A\) 有两张特别的态射
\[
 \alpha_A:\quad 0\longleftarrow0\longrightarrow A,
 \qquad
 \beta_A:\quad 0\longleftarrow A\xrightarrow{1_A}A.
\]
它们在分类空间中给出两条从 \(0\) 到 \(A\) 的路径。以下采用先沿 \(\beta_A\) 前进、再沿 \(\alpha_A\) 反向返回的环路方向，对应 \(A\) 的正可加类；反过来走给出它的负类。对象的可加信息因而可以保留在基本群中，这正是下面将 \(K_0\) 放在 \(\pi_1\) 的原因。

\begin{definition}\label{b4:def:B-higher-k}
设 \(\mathcal E\) 为小正合范畴，以零对象作为 \(BQ\mathcal E\) 的基点。对 \(n\ge0\)，定义
\[
K_n(\mathcal E)=\pi_{n+1}(BQ\mathcal E,0).
\]
对本质小范畴先选取小骨架；等价范畴的神经同伦等价保证此选择不影响所得群。对含幺环 \(R\)，以 \(R\text{-}\mathrm{proj}\) 表示有限生成投射左 \(R\)-模范畴，取分裂正合结构，并定义 \(K_n(R)=K_n(R\text{-}\mathrm{proj})\)。
\end{definition}
\symidx{\(K_n(\mathcal E),K_n(R)\)：正合范畴与环的代数 \(K\) 群}

这些定义与前文的低阶定义相容：
\(\pi_1(BQ\mathcal E)\simeq K_0(\mathcal E)\)，且环上的 \(K_1\) 同构于稳定一般线性群的交换化。由此，前一小节的矩阵计算也给出了本节定义的 \(K_1\)。比较结果见 Weibel\cite[Chapter IV, Proposition 6.2; Section 7]{Weibel2013KBook}。
前一种比较直接计算 \(BQ\mathcal E\) 的基本群；后一种比较使用其环路空间的基点连通分支，将 \(\pi_2(BQ\mathcal E)\) 转成这一分支的基本群，再以加号构造计算它。群完成和加号构造的空间比较仍需所引的拓扑论证。
\begin{proofsketch}
范畴神经的一个单形就是一串可复合态射，因此其基本群可以用边和二单形的复合关系计算。取上述 \(\alpha_A\) 为连接各顶点的树，\(\beta_A\) 后接 \(\alpha_A\) 反向路径的环路类记为 \([A]\)。每个可容许短正合列 \(0\to A\to B\to C\to0\) 给出关系 \([B]=[A]+[C]\)。反过来，每条边都表示一个可容许子商；若由 \(A\leftarrow M\to B\) 表示，其相对于上述树的环路类由 \(\ker(M\to A)\) 的类给出。特别地，\(\alpha_B\) 是树边，给出零类；\(\beta_B\) 给出 \([B]\)。拉回计算说明，复合两条边时相应核的类相加。二单形关系进一步给出可容许短正合列的加法关系；从全部子商关系归约到这些关系的图式追踪见所引低阶比较的证明。直和的两个次序同构，使生成元彼此交换，得到第一项。

对 \(\mathcal E=R\text{-}\mathrm{proj}\)，全部正合列分裂。令 \(\mathcal S\) 为投射模及同构组成的范畴，以直和为运算；再用短正合列组成辅助范畴，按末项投影到 \(Q\mathcal E\)。在两边同时添加一个直和因子，并把此操作形式地变为可逆，可以比较它的纤维与 \(\mathcal S\) 的群完成。拉回短正合列给出纤维之间的比较；分裂性使这些比较为同伦等价。因此群完成空间与 \(\Omega BQ\mathcal E\) 同伦等价。这一步的范畴定义与纤维比较见\cite[Chapter IV, Definition 7.3; Example 7.4; Lemma 7.5; Proposition 7.6; Theorem 7.8]{Weibel2013KBook}，这里省略神经实现与群完成的技术验证。

每个有限生成投射左模都有补模，添加补模后成为有限自由左模。把自由左模 \({}_RR^n\) 的元素写成行向量，则每个自同构唯一写成 \(r_A:x\mapsto xA\)，其中 \(A\in\operatorname{GL}_n(R)\)。由于 \(r_A\circ r_B=r_{BA}\)，其自同构群是 \(\operatorname{GL}_n(R)^{\mathrm{op}}\)。逆矩阵映射 \(A\mapsto A^{-1}\) 给出从这个反群到 \(\operatorname{GL}_n(R)\) 的群同构，与块稳定化及块直和相容，并把初等矩阵 \(e_{ij}(a)\) 送到 \(e_{ij}(-a)\)。因此群完成与加号构造的比较给出
\[
\begin{aligned}
 \pi_0\bigl(\Omega BQ\mathcal E\bigr)&\cong K_0(R),\\
 \bigl(\Omega BQ\mathcal E\bigr)_0
 &\simeq B\bigl(\operatorname{GL}(R)^{\mathrm{op}}\bigr)^+
 \simeq B\operatorname{GL}(R)^+.
\end{aligned}
\]
其中下标 \(0\) 表示基点所在的连通分支。加号构造杀掉初等子群 \(E(R)\)，并保持同调；\cref{b4:lem:B-whitehead}给出 \(E(R)=[\operatorname{GL}(R),\operatorname{GL}(R)]\)，其证明中的初等交换子公式还表明 \(E(R)\) 自身是完美群。

逆矩阵识别把自同构 \(r_A\) 的环路类送到 \([A^{-1}]\)。群完成的零分支是由直和运算给出的群状 \(H\)-空间，它的同伦逆在 \(\pi_1\) 上再次取逆。将上述比较与这项同伦逆复合，便得到把 \([r_A]\) 送到 \([A]\) 的比较；以下使用这一约定。于是
\[
 \pi_2(BQ\mathcal E)\cong\pi_1\bigl((\Omega BQ\mathcal E)_0\bigr)
 \cong\operatorname{GL}(R)/E(R)=K_1(R).
\]
群完成与加号构造的空间比较见\cite[Chapter IV, Theorems 4.10 and 7.1; Corollary 7.2; Exercise 7.9]{Weibel2013KBook}。该书的环版本采用右模；这里先把对称幺半范畴的群完成及分裂正合范畴的一般比较用于左模，再通过反群的逆矩阵识别得到同一加号空间。
\end{proofsketch}

\subsection{局部化与高阶信息}
\label{b4:subsec:B03:03}

高阶 \(K\) 群使零阶的可加关系延伸为长正合列。以下给出一个有明确适用范围的代表性结果。

\ProseParagraphBreak
先说明式中商范畴的意义。交换范畴 \(\mathcal A\) 的 Serre 子范畴 \(\mathcal B\) 对子对象、商对象与扩张封闭。Serre 商 \(\mathcal A/\mathcal B\) 保留原来的对象，并把核、余核都属于 \(\mathcal B\) 的态射变成同构；其中 \(A\) 成为零对象，当且仅当 \(A\in\mathcal B\)。于是包含函子记录被消去的部分，商函子记录消去后仍可见的部分，长列中的三组 \(K\) 群分别属于这三个正合范畴。
\symidx{\(\mathcal A/\mathcal B\)：交换范畴关于 Serre 子范畴的商}

\ProseParagraphBreak
例如在有限生成交换群范畴中消去有限群，有限指数的包含 \(n\mathbb Z\hookrightarrow\mathbb Z\) 的核为零、余核有限，因而在商范畴中可逆。屋顶
\[
 \mathbb Z\longleftarrow n\mathbb Z\longrightarrow\mathbb Z,
 \qquad nm\longmapsto m\quad(n\ge1),
\]
便表示乘以 \(1/n\)。这解释了为何该商范畴的态射会出现有理系数，而不是仅把每个对象的有限挠子群删去后仍保留原有态射。

\begin{theorem}[Quillen 局部化]\label{b4:thm:B-quillen-localization}
设 \(\mathcal A\) 为小交换范畴，\(\mathcal B\subset\mathcal A\) 为 Serre 子范畴，即对子对象、商对象与扩张封闭的全子范畴。则包含与 Serre 商函子给出长正合列，其中中间段的 \(n\ge1\)：
\[
\begin{aligned}
\cdots&\longrightarrow K_{n+1}(\mathcal A/\mathcal B)
 \xrightarrow{\partial}K_n(\mathcal B)
 \longrightarrow K_n(\mathcal A)\\
&\longrightarrow K_n(\mathcal A/\mathcal B)
 \xrightarrow{\partial}K_{n-1}(\mathcal B)\longrightarrow\cdots\\
&\longrightarrow K_0(\mathcal B)\longrightarrow K_0(\mathcal A)
 \longrightarrow K_0(\mathcal A/\mathcal B)\longrightarrow0.
\end{aligned}
\]
本质小情形通过小骨架作同样解释。
\end{theorem}
\begin{proofsketch}
写 \(q:\mathcal A\to\mathcal A/\mathcal B\)，考虑
\(Qq:Q\mathcal A\to Q(\mathcal A/\mathcal B)\)。对商范畴的对象 \(L\)，令 \(\mathcal D_L\) 的对象为 \(A\in\mathcal A\) 连同 \(Q(\mathcal A/\mathcal B)\) 中的态射 \(L\to q(A)\)。这是计算该函子同伦纤维的逗号范畴。其全子范畴 \(\mathcal F_L\) 只保留所给态射为同构的对象。

商范畴中的子商可以由原范畴中的层 \(A_1\subseteq A_2\subseteq A\) 表示。同一个子商的两个表示可通过对子对象取交与和得到共同细化；比较这些表示的映射，其核与余核属于 \(\mathcal B\)，Serre 封闭性使细化中相应的差商仍属于 \(\mathcal B\)。固定一个子商表示后，所有细化所成的范畴是滤过的，因而其神经可缩。Quillen 的定理 A 遂给出 \(B\mathcal F_L\simeq B\mathcal D_L\)。当 \(L=0\) 时，\(q(A)\cong0\) 恰好等价于 \(A\in\mathcal B\)，故
\[
 B\mathcal D_0\simeq B\mathcal F_0\simeq BQ\mathcal B.
\]

还须证明商范畴中的态射引起 \(\mathcal D_L\) 之间的同伦等价。为此，对 \(N\in\mathcal A\) 考察所有在商范畴中成为同构的映射 \(A\to N\)。先取像，便把它化成满射；核属于 \(\mathcal B\)。将核作可容许扩张并用拉回、推出改变表示，可以比较这个辅助范畴与 \(Q\mathcal B\)。所有像 \(N'\subseteq N\)、\(N/N'\in\mathcal B\) 有共同细化，因而取不同的像不改变其同伦型。由此得到所需的换基同伦等价。

Quillen 的定理 B 现在识别出同伦纤维列
\[
 BQ\mathcal B\longrightarrow BQ\mathcal A
 \longrightarrow BQ(\mathcal A/\mathcal B).
\]
同伦群长正合列的典型片段为
\[
 \pi_{n+1}(BQ\mathcal A)
 \longrightarrow\pi_{n+1}\bigl(BQ(\mathcal A/\mathcal B)\bigr)
 \xrightarrow{\partial}\pi_n(BQ\mathcal B),
 \qquad n\ge1.
\]
代入 \(K_n=\pi_{n+1}BQ\)，右端便是 \(K_{n-1}(\mathcal B)\)：连接映射降低一次 \(K\) 次数，来自同伦长正合列的次数变化。由于 \(BQ\mathcal B\) 连通，其 \(\pi_0\) 为零，故末尾的 \(K_0(\mathcal A)\to K_0(\mathcal A/\mathcal B)\) 满射。概要省略定理 A、B 的拓扑证明和上述辅助范畴的全部自然变换；前者见\cite[Chapter IV, Theorems 3.7--3.8]{Weibel2013KBook}，后者按\cite[Chapter V, Theorem 5.1, Claims 5.1.2--5.1.7]{Weibel2013KBook}逐项验证。
\end{proofsketch}

令 \(\mathcal A\) 为有限生成交换群，\(\mathcal B\) 为有限交换群。张量 \(\mathbb Q\) 杀掉且仅杀掉有限群，Serre 商等价于有限维 \(\mathbb Q\)-向量空间范畴。态射上，这个识别来自清除分母：有理线性映射在某个有限指数子群上成为整数线性映射，而有限指数的变化在商范畴中成为同构。有限交换群的合成因子都是素阶循环群，所以短正合关系把每个 \([M]\) 写成 \([\mathbb Z/p\mathbb Z]\) 的有限整数和。另一方面，短正合列中的群阶相乘，故各 \(p\)-主分量的长度 \(v_p(|M|)\) 诱导群同态
\[
K_0(\mathcal B)\longrightarrow\bigoplus_{p\;\text{为素数}}\mathbb Z,
\qquad [M]\longmapsto\bigl(v_p(|M|)\bigr)_p.
\]
它把 \([\mathbb Z/p\mathbb Z]\) 送到第 \(p\) 个基向量，与将该基向量送回 \([\mathbb Z/p\mathbb Z]\) 的群同态互为逆，因而得到 \(K_0(\mathcal B)\simeq\bigoplus_p\mathbb Z\)，没有额外的生成元关系。在 \(K_0(\mathcal A)\) 中，短正合列
\(0\rightarrow\mathbb Z\xrightarrow{p}\mathbb Z\rightarrow\mathbb Z/p\mathbb Z\rightarrow0\)
使这些生成元全为零。有限生成交换群的结构定理因而给出
\(K_0(\mathcal A)\simeq\mathbb Z\)，由秩识别。

\ProseParagraphBreak
连接映射如何产生有限群的类，可以先看素数 \(p\)。整数映射 \(p:\mathbb Z\to\mathbb Z\) 的核为零、余核为 \(\mathbb Z/p\mathbb Z\)，所以它在 Serre 商中成为同构，且对应 \(\mathbb Q\) 上乘 \(p\) 的自同构。在原范畴中，这个提升未能成为同构的部分正是其余核。按取余核类减核类的符号约定，连接映射给出
\[
 \begin{aligned}
 \partial(p)&=[\mathbb Z/p\mathbb Z],\\
 \partial(p^m)&=m[\mathbb Z/p\mathbb Z],\\
 \partial(p^{-m})&=-m[\mathbb Z/p\mathbb Z]\quad(m\ge1).
 \end{aligned}
\]
后一组等式来自 \(\partial\) 从乘法群到加法群的同态性质。乘 \(-1\) 的整数映射本来就是自同构，核与余核都为零，故 \(\partial(-1)=0\)。该余核减核的公式见 Weibel\cite[Chapter V, Application 6.1 and Example 6.1.2]{Weibel2013KBook}。将任意非零有理数分解为符号与有限多个素数幂，就得到素数赋值
\[
\mathbb Q^\times\longrightarrow\bigoplus_p\mathbb Z,
\qquad a\longmapsto\bigl(v_p(a)\bigr)_p.
\]
仅从算术也能直接验证
\[
1\longrightarrow\{1,-1\}\longrightarrow\mathbb Q^\times
\longrightarrow\bigoplus_p\mathbb Z\longrightarrow0
\]
正合：素因子分解给出核与满射性。这说明在零阶可加关系中消失的有限群类，会作为高一阶连接映射的值重新出现。

\begin{proposition}\label{b4:prop:B-stabilization-example}
在域 \(\mathbb F_2\) 上，\(\operatorname{GL}_2(\mathbb F_2)\simeq S_3\)，它的交换化为 \(C_2\)，但稳定交换化 \(K_1(\mathbb F_2)\) 是平凡群。二阶初等矩阵 \(e_{12}(1)\) 在前一个交换化中给出非平凡类，添加一个坐标后成为交换子。
\end{proposition}
\begin{proof}
可逆矩阵置换 \(\mathbb F_2^2\) 的三个非零向量。如果某矩阵固定它们，就特别固定一组基，因而是恒等矩阵，所以这个作用忠实。可逆矩阵的第一列有 \(3\) 种选择，第二列有 \(2\) 种选择，总数为 \(6\)，故所得到 \(S_3\) 的单同态是同构。\(S_3\) 的符号同态核为 \(A_3\)，而两个换位的交换子给出一个三轮换，故其交换子群恰为 \(A_3\)，交换化为 \(C_2\)。矩阵 \(e_{12}(1)\) 固定 \((1,0)^t\)，交换 \((0,1)^t\) 与 \((1,1)^t\)，所以对应换位并给出非平凡交换化类。稳定化后有
\[
 e_{12}(1)=\bigl[e_{13}(1),e_{32}(1)\bigr],
\]
故这个类在稳定交换化中变为单位类。最后，\cref{b4:thm:B-k1-field-integers}给出 \(K_1(\mathbb F_2)\simeq\mathbb F_2^\times=\{1\}\)。
\end{proof}

\begin{proposition}\label{b4:prop:B-q-zero-subspaces}
设 \(q\) 为素数幂，\(\mathcal E\) 为以 \(\mathbb F_q^n\)（\(n\ge0\)）为对象的有限维向量空间范畴的标准小骨架，取通常的正合结构。对 \(V\in\mathcal E\)，\(Q\mathcal E\) 中从 \(0\) 到 \(V\) 的态射与 \(V\) 的子空间一一对应。特别地，
\[
 \#\operatorname{Hom}_{Q\mathcal E}(0,\mathbb F_q^2)=q+3.
\]
原范畴的零对象 \(0\) 在 \(Q\mathcal E\) 中既不是始对象，也不是终对象。
\end{proposition}
\begin{proof}
从 \(0\) 到 \(V\) 的态射由 \(0\leftarrow M\xrightarrow{i}V\) 表示。到零对象的映射自动满射，而单射 \(i\) 把 \(M\) 识别为 \(V\) 中的子空间 \(i(M)\)。两张图式等价当且仅当它们有相同的像：若像相同，经各自单射识别便得到使图式相容的中间对象同构；若图式等价，则像当然相同。每个子空间又能选取骨架中的同构代表，故得到所述对应。

二维空间的零子空间与全空间各一个；一维子空间的非零向量彼此不交，每个含 \(q-1\) 个非零向量，故一维子空间有 \((q^2-1)/(q-1)=q+1\) 个。总数为 \(q+3\)，所以 \(0\) 不是始对象。若 \(V\ne0\)，一张 \(V\to0\) 的 \(Q\)-态射需要单射 \(M\to0\)，迫使 \(M=0\)，却又要求 \(M\to V\) 满射，矛盾。因此 \(0\) 也不是终对象。
\end{proof}
